\section{Architecture and Interfaces}

\subsection{Top-Level Architecture}
\begin{figure}[H]
    \centering
    \begin{tikzpicture}
        \draw[brandgreen, line width=1.1pt, rounded corners=3pt] (0,0) rectangle (11,5.6);
        \draw[brandgreen, line width=1pt] (0.8,1.8) rectangle (3.0,4.3);
        \draw[brandgreen, line width=1pt] (4.4,1.8) rectangle (6.6,4.3);
        \draw[brandgreen, line width=1pt] (8.0,1.8) rectangle (10.2,4.3);
        \draw[brandgreen, -latex, line width=1pt] (3.0,3.05) -- (4.4,3.05);
        \draw[brandgreen, -latex, line width=1pt] (6.6,3.05) -- (8.0,3.05);
        \node at (1.9,3.05) {Input};
        \node at (5.5,3.05) {Processing};
        \node at (9.1,3.05) {Output};
        \node[text width=8.8cm, align=center] at (5.5,0.7) {Placeholder for system block diagram, control architecture, or subsystem partitioning};
    \end{tikzpicture}
    \caption{Architecture overview}
    \label{fig:architecture}
\end{figure}

\subsection{Interface Summary}
\begin{table}[H]
    \centering
    \small
    \begin{tabular}{p{2.5cm} p{2.3cm} p{2.1cm} p{4.0cm}}
        \toprule
        \textbf{Interface} & \textbf{Direction} & \textbf{Domain} & \textbf{Notes} \\
        \midrule
        Power In & Input & Electrical & Define range, protection, and startup expectations. \\
        Control Bus & Bidirectional & Digital & Include timing, protocol, and fault handling notes. \\
        Sensor Path & Input & Analog / Mixed & State filtering, scaling, and calibration assumptions. \\
        Status Output & Output & Digital / UI & Define fault encoding and update cadence. \\
        \bottomrule
    \end{tabular}
    \caption{Example interface matrix}
    \label{tab:interfaces}
\end{table}

\decisionbox{
Document the architecture decision directly: which partitioning was selected,
what was rejected, and what assumption still carries the most risk.
}
